% MT041 / explanation — canonical editable content.
二分类LR先得到任意实数分数\(f=w^Tx+b\)，再计算



\[
P(y=1\mid x)=\sigma(f)=\frac{1}{1+e^{-f}}.
\]


因为\(e^{-f}>0\)，结果始终在0与1之间。\(f=0\)时概率为0.5，\(f\)越大越偏向类1；这就是B、D的依据。预测时再按阈值把概率变成类别。



\begin{itemize}

\item \textbf{A错误。} LR直接用上式表示给定\(x\)后的类别概率，没有“类别似然除以输入先验”的计算。Bayes公式也应包含类别先验，选项混用了概念。

\item \textbf{B正确。} Sigmoid把实数分数映射成0到1之间的数。

\item \textbf{C错误。} One-vs-rest（OvR）是“一类对其余所有类”。三类时分别训练1对其余、2对其余、3对其余，共3个二分类器，再比较分数。它不是任意画两条线分三个区域。

\item \textbf{D正确。} 输出按模型定义解释为类别概率；模型假设或拟合不佳时，概率未必在真实数据中校准准确。

\item \textbf{E错误。} 在附近小改参数也不能降低损失，叫局部最小；在全部可能的参数中损失最低，叫全局最小。凸目标的局部最小值也是全局最小值，但最优参数可能不止一组，最低损失也未必能在有限参数处达到。无正则且数据完全可分就是后一种情况；例子见\StudyRef{Answers}{mt001}{MT001}。迭代是否收敛还取决于学习率等设置。

\end{itemize}


\textbf{出处：} 2023A期中 Q2，原卷第2页（5分）；2025A期中 Q2，原卷第2页（5分）


\textbf{答案依据：} 2023A期中 Q2 随卷／配套参考答案；条件与解释经复核。
